% paper/sections/05_advection.tex
%
% §5  CLS 移流の離散化：Dissipative CCD スキーム・時間積分・安定性制約
%     （旧 04_time_integration.tex + 04b_time_schemes.tex の統合）
%     WENO5 は参照スキームとして付録~\ref{app:weno5} に収録する．
%
\section{CLS 移流の離散化}
\label{sec:time}

本章では，第~\ref{sec:levelset}章で導入した CLS 移流方程式
\begin{equation}
  \frac{\partial\psi}{\partial t} + \nabla\cdot(\psi\bm{u}) = 0
  \label{eq:cls_advection_ref}
\end{equation}
（式~\eqref{eq:cls_advection}の再掲）を数値的に解くための離散化スキームを示す．
第~\ref{sec:CCD}章で構築した CCD 法および第~\ref{sec:dissipative_ccd}節の
Dissipative CCD フィルタを移流方程式へ適用する設計根拠を述べ，
圧縮項（再初期化）・時間積分・安定性制約と合わせて全スキームを定式化する．

% ─────────────────────────────────────────────────────────────────────────────
\subsection{Dissipative CCD による CLS 移流の安定化：設計根拠}
\label{sec:advection_motivation}
% ─────────────────────────────────────────────────────────────────────────────

% ─────────────────────────────────────────────────────────────────────────────
\subsubsection{移流スキーム選択の設計根拠：WENO5 からの脱却}
\label{sec:advection_weno5_critique}
% ─────────────────────────────────────────────────────────────────────────────

\paragraph{WENO5 の設計思想と CLS 界面の非適合性．}
WENO5（Weighted Essentially Non-Oscillatory 5次精度）スキーム~\cite{JiangShu1996}は，
衝撃波・接触不連続面を含む圧縮性流体力学における数値振動（Gibbs 現象）を
強力に抑制するために設計された．
ステンシル内の\textbf{滑らかさ指標}（$\beta_k$，式~\eqref{eq:weno5_beta}）が
不連続を検知すると，その方向の重みを $O(\varepsilon^2)$ まで低減して
\textbf{非線形散逸}を自動的に投入する機構がスキームの中核にある．

\medskip

一方，本稿で採用する保存型レベルセット（CLS）法では，界面を滑らかな
正則化ヘヴィサイド関数 $H_\varepsilon(\phi)$（式~\eqref{eq:Heps_def_preview}）で表現し，
界面厚みを $2\varepsilon \sim O(h)$ に制御する．
この枠組みにおいて $\psi = H_\varepsilon(\phi)$ は\textbf{実質的に滑らかな関数}であり，
移流方程式の解の中に\textbf{真の不連続は存在しない}．
結果として WENO5 の滑らかさ指標は，
滑らかな界面プロファイルを「不連続の兆候」と過剰に検知し，
\textbf{不必要な非線形散逸を常時付加する}こととなる．

この過剰散逸の帰結は以下の 2 点に集約される：
\begin{enumerate}
  \item \textbf{界面の非物理的平滑化：}
        液滴の微細変形，薄液膜の形成，界面の強いせん断巻き込みといった
        物理的に意味のある界面構造が，数値散逸によって均されてしまう．
        特に高 $\mathrm{We}$ 数・高 $\mathrm{Re}$ 数流れで顕著となる．
  \item \textbf{Balanced-Force 精度の侵食：}
        第~\ref{sec:balanced_force}章で導く Balanced-Force 条件は，
        曲率 $\kappa$ の計算と圧力勾配 $\nabla p$ に\textbf{同一の差分演算子}を
        用いることで寄生流れを $O(h^6)$ まで低減する．
        WENO5 の非線形重みは $\psi$ の移流精度を適応的に劣化させるため，
        再初期化後の $\psi$ プロファイル（曲率計算の入力）に振動・ぼやけを残し，
        Balanced-Force によって確保したはずの高精度を実用上無意味にする危険がある．
\end{enumerate}

\paragraph{CCD との根本的親和性．}
CCD（結合コンパクト差分，第~\ref{sec:CCD}章）は，
スペクトル解像度に優れる\textbf{線形・コンパクト・低散逸}スキームである．
CLS が提供する滑らかな $\psi$ プロファイルは，まさに CCD が最も得意とする対象であり，
以下の理由から移流項の離散化として理論的に最適な選択となる．

\begin{description}
  \item[\upshape(i) 設定界面構造の長時間保持]
        非線形重み調整をもたない CCD は，
        設定した $\varepsilon$ のプロファイルを移流中に追加散逸なく維持する．
        WENO5 が数ステップで界面を $O(h)$ 幅だけ肥厚させる状況でも，
        CCD は高忠実な界面形状を保ち続ける（第~\ref{sec:verification}章参照）．
  \item[\upshape(ii) 空間離散化スキームの統一]
        本フレームワークでは PPE（§\ref{sec:pressure}），粘性項，
        表面張力の曲率評価，速度場微分の全てに CCD を用いている．
        移流項に\textbf{同一の 3 点コンパクトステンシル}を採用することで，
        支配方程式全体の空間離散化が単一の構造で統一される．
        この統一性は，境界条件の処理を一元化し（§\ref{sec:ccd_bc}），
        変数間の数値的干渉を低減する——
        WENO5 の 5 点ステンシルと CCD の 3 点ステンシルを混在させた場合に生じる
        \textbf{打ち切り誤差の非整合}（$O(h^5)$ の移流誤差が $O(h^6)$ の圧力精度を
        律速する非均質な精度階層）を回避できる．
\end{description}

% ─────────────────────────────────────────────────────────────────────────────
\subsubsection{気液界面移流の困難と純 CCD の限界}
\label{sec:advection_ccd_limit}

$\psi$ は界面付近で $0$ から $1$ へと急激に変化する（界面幅 $\sim 2\varepsilon$）．
この急峻なプロファイルを移流方程式~\eqref{eq:cls_advection_ref}で輸送するとき，
\textbf{中心差分ベースのスキームは原理的に不安定}になる．

フォン・ノイマン安定性解析による確認：
CCD は純中心差分（修正波数 $k^*(kh)$ が実数）であるため，
移流方程式 $\partial\psi/\partial t + u\,\partial\psi/\partial x = 0$ に
陽的時間積分を組み合わせると，Fourier モード $\psi_i = e^{i\xi i}$ に対して
増幅係数は
\begin{equation}
  |g(\xi)|^2 = 1 + \sigma^2 \left[\frac{k^*(\xi)}{k}\right]^2 > 1
  \quad \forall\, \xi > 0, \;\sigma = \frac{u\Delta t}{h} > 0
  \label{eq:ccd_adv_instability}
\end{equation}
となり，\textbf{すべての波数で振幅が増大する}——これは純 CCD による陽的移流が
本質的に不安定であることを意味する（CFL 条件を小さくしても解消しない）．

この不安定性の物理的帰結は \textbf{Gibbs 振動}（$\psi < 0$ や $\psi > 1$
の非物理的値の発生）であり，界面のにじみおよび密度場の不整合として現れる
（§\ref{sec:failure_modes}）．

\subsubsection{Dissipative CCD の採用：フィルタによるスペクトル安定化}
\label{sec:advection_dccd_design}

第~\ref{sec:dissipative_ccd}節で設計した 3 点スペクトルフィルタ
（式~\eqref{eq:dccd_filter}）を，
CLS 移流フラックス発散の計算に組み込むことで上記の不安定性を解消する．

\medskip
\noindent\textbf{移流への適用．}
フラックス場 $f_i = \psi_i u_i$（$x$ 方向）に CCD を適用して微分場 $f'_i$
を求め，その後 Dissipative フィルタを適用する：
\begin{align}
  f'_i &= \left[\frac{\partial(\psi u)}{\partial x}\right]_i
  \quad\text{（CCD 求解，$O(h^6)$ 精度）}
  \label{eq:dccd_adv_ccd} \\[4pt]
  \tilde{F}_i &= f'_i + \varepsilon_d^{\mathrm{adv}}\!\left(f'_{i+1} - 2f'_i + f'_{i-1}\right)
  \label{eq:dccd_adv_filter}
\end{align}
得られた修正フラックス発散 $\tilde{F}_i$ を時間積分（§\ref{sec:time_int}）の
空間演算子として使用する：$\mathcal{L}(\psi) = \tilde{F}$．

\medskip
\noindent\textbf{安定化機構：フィルタ伝達関数による振幅減衰．}
Fourier モード $\psi_i = e^{i\xi i}$ に対して
$f'_{i+1} - 2f'_i + f'_{i-1} = (e^{i\xi}-2+e^{-i\xi})f'_i = -4\sin^2(\xi/2)\,f'_i$
であるから，フィルタ適用後の有効スペクトルは
$\tilde{F}(\xi) = H(\xi;\varepsilon_d^{\mathrm{adv}})\cdot f'(\xi)$，
$H(\xi) = 1 - 4\varepsilon_d^{\mathrm{adv}}\sin^2(\xi/2) \leq 1$．

この実数係数（$0 \leq H \leq 1$）の乗算が\textbf{振幅を単調に減衰}させ，
増幅係数を
\begin{equation}
  |g(\xi)|^2 = 1 + \sigma^2 H(\xi)^2\left[\frac{k^*(\xi)}{k}\right]^2
  \label{eq:dccd_adv_amp}
\end{equation}
に変える．$H(\xi) < 1$ によって有効 CFL 数が $\sigma H(\xi) < \sigma$ に縮小し，
高波数（$\xi \to \pi$，$H \to 1-4\varepsilon_d^{\mathrm{adv}}$）では
成長率が大幅に低減される．
典型的な二相流パラメータ（$\mathrm{CFL} \leq 0.5$，$\varepsilon_d^{\mathrm{adv}} = 0.05$）のもとで
実用的な安定性が得られることを数値実験で確認している（第~\ref{sec:verification}章参照）．

\begin{tcolorbox}[warnbox, title={移流フィルタの散逸強度：S(ψ) 制御を使わない理由}]
\label{warn:adv_filter_uniform}

第~\ref{sec:dissipative_ccd}節の速度・圧力場フィルタでは
スイッチ関数 $S(\psi) = (2\psi-1)^2$ を用いて
\textbf{界面でフィルタを OFF}（$S(1/2) = 0$）にし，曲率計算の精度を保護した．

CLS 移流に同じ $S(\psi)$ を適用すると，Gibbs 振動が最も激しい
$\psi \approx 0.5$（界面）の近傍でフィルタ強度がゼロになり，
安定化が最も必要な場所で無効化されるという\textbf{設計矛盾}が生じる．

\textbf{したがって CLS 移流には一様散逸強度を使用する：}
\begin{equation}
  \varepsilon_d^{\mathrm{adv}} = \varepsilon_{d,\max} = 0.05 \quad\text{（全格子点，S(ψ) 制御なし）}
  \label{eq:eps_adv}
\end{equation}
この値のもとでナイキスト波長 $2h$（$\xi=\pi$）での伝達関数は
$H(\pi;\,0.05) = 1 - 4\times 0.05 = 0.80$（20\%\ 振幅減衰）；
物理的に重要な低波数成分（$\xi \ll \pi$）への影響は
$H(\xi) \approx 1 - \varepsilon_d^{\mathrm{adv}}\xi^2 \approx 1$（精度劣化 $O(\varepsilon_d\xi^2) \ll 1$）に抑えられる
（式~\eqref{eq:dccd_transfer}，§\ref{sec:dccd_filter_theory}）．
典型的な CFL $\leq 0.5$ のもとでの実用的安定性は第~\ref{sec:verification}章のベンチマークで検証する．

速度・圧力場フィルタ（適応制御 $\varepsilon_d^{(i)} = \varepsilon_{d,\max}S(\psi_i)$）
との使い分けを下表に整理する．
\end{tcolorbox}

\medskip
\noindent\textbf{質量保存性．}
フィルタ項 $\varepsilon_d^{\mathrm{adv}}(f'_{i+1}-2f'_i+f'_{i-1})$ の周期境界での総和は
\[
  \sum_i \varepsilon_d^{\mathrm{adv}}\!\left(f'_{i+1} - 2f'_i + f'_{i-1}\right)
  = \varepsilon_d^{\mathrm{adv}}\!\left(\sum_i f'_{i+1} - 2\sum_i f'_i + \sum_i f'_{i-1}\right) = 0
\]
（添字のシフト対称性より）．
すなわち\textbf{フィルタは総質量に寄与しない}（周期 BC）．
CCD そのものは $\sum_i f'_i\,h \approx [f]_0^L$（発散定理）を $O(h^6)$ 精度で近似するため，
$N$ 格子点での局所打ち切り誤差 $O(h^6/h)=O(h^5)$ が $N$ 点累積して
$h \cdot N \cdot O(h^5) = O(h^5 \cdot L)$ となり，
移流ステップあたりの質量保存誤差は $O(h^5\,\Delta t)$ となる
（WENO5 のフラックス差分法による完全な保存性と比較する場合は付録~\ref{app:weno5} 参照）．

\begin{tcolorbox}[warnbox, title={TVD 性の欠如と ψ のクランプ処置}]
\label{warn:adv_clamp}

WENO5 + TVD-RK3 は TVD 性（Total Variation Diminishing）により
$\psi \in [0,1]$ を自動保証する．
Dissipative CCD は線形フィルタであり TVD 性を持たないため，
界面近傍での Gibbs 振動が抑制不足になると $\psi$ が $[0,1]$ を
逸脱する可能性がある．

\textbf{対処：} 各 TVD-RK3 ステージ完了後に
\begin{equation}
  \psi_i \leftarrow \max\bigl(0,\, \min(1,\, \psi_i)\bigr)
  \label{eq:psi_clamp}
\end{equation}
のクランプを適用する．クランプによる質量誤差は $O(h^2)$ で
CSF モデル誤差 $O(\varepsilon^2) \approx O(h^2)$（§\ref{sec:balanced_force}）
と同オーダーであり，全体精度のボトルネックにならない．
クランプ後の質量変化は第~\ref{sec:verification}章の体積保存ベンチマークで定量確認すること．
\end{tcolorbox}

\begin{tcolorbox}[warnbox, title={線形 CCD 移流の残余リスクと既知の失敗モード}]
\label{warn:adv_risks}

Dissipative CCD とクランプ処置（式~\eqref{eq:psi_clamp}）は
TVD 欠如リスクを実用的に抑制するが，
以下の 2 点はスキームの\textbf{設計上の固有制限}として明示しておく．

\medskip
\noindent\textbf{(A) 非線形エイリアシングの蓄積．}
線形スキームである CCD は振幅の自律的減衰機能を持たない．
移流項 $\nabla\cdot(\psi\bm{u})$ の非線形相互作用（$\psi$ と $\bm{u}$ の積）は，
低波数 $\to$ 高波数へのエネルギーカスケードを生成し，
格子スケールのノイズとして蓄積する．
本フレームワークではこれに対して二重の対策を施している：
\begin{itemize}
  \item Dissipative CCD フィルタ（$\varepsilon_d^{\mathrm{adv}} = 0.05$）が
        高波数成分（$\xi \to \pi$）を各ステップで 20\%\ 減衰させる
        （$H(\pi;\,0.05) = 0.80$，式~\eqref{eq:eps_adv}）；
  \item Rhie--Chow 補間（§\ref{sec:rhie_chow}）が
        コロケート格子上でのチェッカーボード励起を独立に抑制する．
\end{itemize}
高 $\mathrm{Re}$ 数・高解像度（$N \geq 256$）ケースにおいて
これらの対策が十分かどうかは，第~\ref{sec:verification}章の
乱流混合ベンチマークで評価すること
（現時点では層流・中低 $\mathrm{Re}$ 数の検証にとどまる）．

\medskip
\noindent\textbf{(B) 極薄界面設定時の破綻．}
界面厚みを格子幅に極限まで近づけた設定
（$\varepsilon \to h$，すなわち $\varepsilon_{\mathrm{factor}} \to 1.0$）では，
$\psi$ プロファイルがほぼ step 関数に退化し，
CCD の高精度近似前提（関数の十分な滑らかさ）が成立しなくなる．
この条件下では Dissipative CCD の散逸補正が不十分となり，
激しい振動とクランプによる質量損失が急増する．

\textbf{安全な運用範囲：} $\varepsilon_{\mathrm{factor}} \geq 1.5$
（デフォルト値，\texttt{NumericsConfig.epsilon\_factor}）を推奨する．
$\varepsilon_{\mathrm{factor}} < 1.2$ かつ流れが非線形に強い場合
（$\mathrm{We} > 100$，密度比 $> 100$）は
計算破綻のリスクがあり，WENO5（付録~\ref{app:weno5}）への
フォールバックを検討すること．
\end{tcolorbox}

\medskip
\noindent\textbf{アルゴリズムの概要（1 次元，TVD-RK3 第 1 ステージ）：}

\begin{tcolorbox}[algbox, title={Dissipative CCD による CLS 移流アルゴリズム（1D)}]
\label{alg:dccd_adv}

\textbf{入力：} $\{\psi_i^n\}$，$\{u_i^n\}$，散逸強度 $\varepsilon_d^{\mathrm{adv}}$

\begin{enumerate}
  \item \textbf{フラックス場の構成：}
        各格子点 $i$ で $f_i = \psi_i \cdot u_i$．

  \item \textbf{CCD 求解（Eq-I + Eq-II システム）：}
        §\ref{sec:ccd_matrix} のブロック Thomas 法で
        $\{f'_i\} = \{(\partial(\psi u)/\partial x)_i\}$ を得る．

  \item \textbf{Dissipative フィルタの適用（一様 $\varepsilon_d^{\mathrm{adv}}$）：}
        内点 $i = 1,\ldots,N-2$：
        \[
          \tilde{F}_i = f'_i + \varepsilon_d^{\mathrm{adv}}\!\left(f'_{i+1} - 2f'_i + f'_{i-1}\right)
        \]
        周期 BC：折り返しインデックスで $i=0,N-1$ にも同様に適用．
        壁面 BC：境界点はフィルタ前の $f'$ 値を使用（§\ref{sec:dccd_bc}）．

  \item \textbf{TVD-RK3 第 1 ステージ更新：}
        \[
          \psi_i^{(1)} = \psi_i^n - \Delta t\,\tilde{F}_i
        \]

  \item \textbf{クランプ：} $\psi_i^{(1)} = \max(0,\min(1,\psi_i^{(1)}))$ （式~\eqref{eq:psi_clamp}）．

  \item 第 2・第 3 RK ステージでステップ 1--5 を繰り返す（式~\eqref{eq:tvd_rk3}）．
\end{enumerate}

\textbf{2 次元への拡張：} 次元分割法（dimensional splitting）を適用し，
$x$ 方向フラックス $\psi u$ と $y$ 方向フラックス $\psi v$ に独立に
1 次元手順を実行する：
\[
  \tilde{F}_{x,ij} = \left[\frac{\partial(\psi u)}{\partial x}\right]_{ij}^{\text{filtered}},
  \quad
  \tilde{F}_{y,ij} = \left[\frac{\partial(\psi v)}{\partial y}\right]_{ij}^{\text{filtered}}
\]
\[
  \psi_{ij}^{(1)} = \psi_{ij}^n - \Delta t\!\left(\tilde{F}_{x,ij} + \tilde{F}_{y,ij}\right)
\]
\end{tcolorbox}

\begin{tcolorbox}[resultbox, title={本稿の空間離散化スキームの役割分担}]
\label{box:scheme_roles}
\begin{center}
\renewcommand{\arraystretch}{1.3}
\begin{tabular}{>{\raggedright\arraybackslash}p{2.5cm}>{\raggedright\arraybackslash}p{3.8cm}>{\raggedright\arraybackslash}p{1.5cm}>{\raggedright\arraybackslash}p{4cm}}
\hline
対象 & スキーム & 精度 & 根拠 \\
\hline
CLS 移流（$\psi$） & Dissipative CCD（§\ref{sec:advection_motivation}） & $O(h^6)$＋フィルタ & 同一ステンシル；一様 $\varepsilon_d^{\mathrm{adv}}$ でGibbs 振動抑制 \\
CLS 圧縮（移流部） & Dissipative CCD（§\ref{sec:cls_compression}） & $O(h^6)$＋フィルタ & 移流と同構成（陽解法） \\
CLS 圧縮（拡散部） & CCD-CN（§\ref{sec:cls_compression}） & $O(h^6)$ & 無条件安定（陰解法） \\
曲率・速度微分 & CCD（§\ref{sec:CCD}） & $O(h^6)$ & 高精度・コンパクト \\
速度・圧力場ノイズ除去 & 散逸的 CCD 適応制御（§\ref{sec:dissipative_ccd}） & — & $S(\psi)$ で界面保護 \\
圧力 Poisson & CCD-Poisson（§\ref{sec:pressure}） & $O(h^6)$ & Balanced-Force 整合性 \\
\hline
\end{tabular}
\end{center}
\end{tcolorbox}

% ─────────────────────────────────────────────────────────────────────────────
\subsection{時間積分スキーム}
\label{sec:time_int}
% ─────────────────────────────────────────────────────────────────────────────

\subsubsection{3次 TVD Runge--Kutta 法（CLS 移流・再初期化のみ）}

\textbf{CLS 移流ステップ（界面追跡）および CLS 再初期化の仮想時間積分}に，
Shu--Osher 型 TVD-RK3 法 \cite{ShuOsher1988} を用いる：
\begin{align}
  q^{(1)} &= q^n + \Delta t \, \mathcal{L}(q^n) \notag\\
  q^{(2)} &= \frac{3}{4}q^n + \frac{1}{4}q^{(1)} + \frac{1}{4}\Delta t\,\mathcal{L}(q^{(1)}) \notag\\
  q^{n+1} &= \frac{1}{3}q^n + \frac{2}{3}q^{(2)} + \frac{2}{3}\Delta t\,\mathcal{L}(q^{(2)})
  \label{eq:tvd_rk3}
\end{align}
ここで $\mathcal{L}(q)$ は空間演算子であり，
CLS 移流では式~\eqref{eq:dccd_adv_filter} の Dissipative CCD フラックス発散
$\mathcal{L}(\psi) = \tilde{F}$（アルゴリズム~\ref{alg:dccd_adv}）を指す．
\textbf{NS 方程式の時間積分には TVD-RK3 は使用しない}（§\ref{warn:tvd_rk3_scope}参照）．

\noindent\textbf{TVD（Total Variation Diminishing）の定義：}
TVD（Total Variation Diminishing）とは，数値解の全変動
$TV(q) \equiv \sum_i |q_{i+1} - q_i|$ が時間的に増大しない性質を指す：
$TV(q^{n+1}) \leq TV(q^n)$．
TVD-RK3 はこの性質を時間方向に保証するが，
Dissipative CCD は線形フィルタであり空間 TVD 性を持たない（§\ref{warn:adv_clamp}）．
クランプ処置（式~\eqref{eq:psi_clamp}）と組み合わせることで
$\psi \in [0,1]$ を維持する．

\begin{tcolorbox}[warnbox, title={TVD-RK3 の適用範囲と全体時間精度}]
\label{warn:tvd_rk3_scope}
\textbf{TVD-RK3 の適用先（CLS のみ）}

TVD-RK3 は以下の2箇所にのみ適用する：
\begin{itemize}
  \item \textbf{CLS 移流方程式}（式 \eqref{eq:cls_advection}）：$\pd{\psi}{t} + \bnabla\cdot(\psi\,\bu) = 0$（保存形）
  \item \textbf{CLS 再初期化方程式}（式 \eqref{eq:cls_reinit_iso}）の仮想時間 $\tau$ 積分（移流部のみ）
\end{itemize}

\medskip
\noindent\textbf{NS 方程式の時間積分方式（AB2 + IPC 法）}

本稿では，プロジェクション法の全体時間精度を $\mathcal{O}(\Delta t^2)$ とするため，
対流項に \textbf{2段 Adams-Bashforth（AB2）法}を採用し，
van Kan（1986）\cite{vanKan1986} の \textbf{Incremental Pressure Correction（IPC）法}と組み合わせる．
Predictor ステップの離散式は：
\begin{equation}
\begin{split}
  \frac{\rho^{n+1}}{\Delta t}(\bu^* - \bu^n)
  &= -\rho^{n+1}\!\left[\tfrac{3}{2}\mathcal{C}(\bu^n,\bu^n)
    - \tfrac{1}{2}\mathcal{C}(\bu^{n-1},\bu^{n-1})\right] \\
  &\quad + \tfrac{1}{2Re}\bigl[\mathcal{D}(\mu^{n+1},\bu^*)
    + \mathcal{D}(\mu^n,\bu^n)\bigr]
  - \bnabla p^n
  + \bm{f}^{n+1}
\end{split}
  \label{eq:predictor_ab2_ipc}
\end{equation}
ここで $\mathcal{C}(\bu,q) = \bu\cdot\bnabla q$（対流演算子），$\mathcal{D}$ は完全粘性応力演算子
（式~\eqref{eq:viscous_full_x}）．
体積力 $\bm{f}^{n+1}$ は表面張力項（$\rho$ 因子なし）と重力項（$\rho^{n+1}$ 因子あり）の和：
\begin{equation*}
  \bm{f}^{n+1} = \frac{\kappa^{n+1}\bnabla\psi^{n+1}}{We}
                - \frac{\rho^{n+1}}{Fr^2}\hat{\bm{y}}
\end{equation*}
（表面張力は CSF 体積力 $\sigma\kappa\bnabla\psi$（式~\eqref{eq:csf_final}）の無次元形；
重力は式~\eqref{eq:gravity_convention} の定義 $\bm{g}^* = -g\hat{\bm{z}}$ に対応し，
2次元鉛直問題では $\hat{\bm{z}} = \hat{\bm{y}}$）．
対流項の外挿係数 $(3/2,\,-1/2)$ が AB2 であり，
前ステップ圧力 $-\bnabla p^n$ の陽的組込みが IPC の本質である．

\medskip
\noindent\textbf{IPC 法の核心：圧力増分 $\delta p$ の求解}%
\label{sec:ipc_derivation}

IPC 法では，圧力 Poisson 方程式において絶対圧力 $p^{n+1}$ の代わりに
\textbf{圧力増分 $\delta p \equiv p^{n+1} - p^n$} を未知変数として解く
（§\ref{sec:algorithm} Step 6--7 参照）：
\begin{align*}
  &\text{PPE：}\quad \bnabla\cdot\!\left(\frac{1}{\rho^{n+1}}\bnabla(\delta p)\right)
    = \frac{1}{\Delta t}\bnabla\cdot\bu^*\\
  &\text{補正：}\quad \bu^{n+1} = \bu^* - \frac{\Delta t}{\rho^{n+1}}\bnabla(\delta p),
  \qquad p^{n+1} = p^n + \delta p
\end{align*}
CCD 演算子 $\mathcal{L}^{\rho}_{\mathrm{CCD}}$ は絶対圧力版と\textbf{完全に同一}であり，
求解対象変数が $p^{n+1}$ から $\delta p$ に変わるだけである．
この定式化により，標準 Chorin 法の人工的圧力境界層に起因する
$\Ord{\Delta t}$ スプリッティング誤差が $\Ord{\Delta t^2}$ に低減される
（Guermond et al., \textit{Acta Numerica}, 2006 \cite{Guermond2006}）．

\medskip
\noindent\textbf{全体的な時間精度}%
\label{sec:time_accuracy_table}

\begin{center}
\renewcommand{\arraystretch}{1.25}
\begin{tabular}{lcc}
\hline
方程式・項 & 時間積分法 & 時間精度 \\
\hline
CLS 移流 & TVD-RK3 & $\Ord{\Delta t^3}$ \\
NS 対流項（Predictor） & AB2（van Kan + 外挿係数 3/2, -1/2） & $\Ord{\Delta t^2}$ \\
NS 粘性項 & Crank--Nicolson & $\Ord{\Delta t^2}$ \\
Projection スプリッティング & IPC（van Kan 1986） & $\Ord{\Delta t^2}$ \\
\hline
\textbf{全体} & & $\boldsymbol{\Ord{\Delta t^2}}$ \\
\hline
\end{tabular}
\end{center}

本スキームの\textbf{全体時間精度は $\Ord{\Delta t^2}$（2次精度）}である．
空間精度（Dissipative CCD 移流：実効 $O(\varepsilon_d^{\mathrm{adv}} h^2)$ フィルタ誤差，
CCD：$O(h^6)$ 微分精度）は，CSF モデル誤差 $O(\varepsilon^2) \approx O(h^2)$
（§\ref{sec:balanced_force}）が律速する全体空間精度と整合する．

\medskip
\noindent\textbf{設計上の注意：「CLS は $\Ord{\Delta t^3}$ なのに全体が $\Ord{\Delta t^2}$ な理由」}

CLS 移流に TVD-RK3（$\Ord{\Delta t^3}$）を採用している一方，NS 対流項は AB2（$\Ord{\Delta t^2}$）にとどめている．
この非対称性は意図的な設計選択であり，以下の理由による：

界面と流体運動は\textbf{疎結合（operator-split）}として扱われる．
CLS 移流（界面輸送）は密度場 $\rho^{n+1}$ と曲率 $\kappa^{n+1}$ を決定する「前処理」であり，
その時間精度が $\Ord{\Delta t^2}$ 以上であれば，全体の流体運動精度は NS 方程式側の $\Ord{\Delta t^2}$ で律速する：
\[
  \text{全体時間精度} = \min\!\bigl(\underbrace{\Ord{\Delta t^3}}_\text{CLS},\;
    \underbrace{\Ord{\Delta t^2}}_\text{NS (AB2+IPC)}\bigr) = \Ord{\Delta t^2}
\]
\textbf{CLS の $\Ord{\Delta t^3}$ 高精度化の役割：}
時間精度を $\Ord{\Delta t^3}$ に向上させることには貢献しないが，
界面プロファイルの変形が小さく保たれることで
（i）曲率 $\kappa$ の数値精度が向上し寄生流れを抑制，
（ii）再初期化の必要頻度が減少し質量保存誤差が低減する，
という実用上の優位が得られる．

\medskip
\noindent\textbf{背景：標準 Chorin 法のスプリッティング誤差（なぜ AB2 だけでは不十分か）}

対流項を AB2（2次精度）に変更するだけでは全体精度は $\Ord{\Delta t^2}$ に\textbf{ならない}．
標準的な Predictor--Corrector 分割（Chorin, 1968）は，
圧力境界条件の不整合（人工的圧力境界層）に起因する $\Ord{\Delta t}$ スプリッティング誤差を独立に内包するため，
対流項の高次化だけで律速が解消されない．
IPC 法による予測ステップへの $p^n$ の組み込みがスプリッティング誤差を除去する唯一の手段である．

\medskip
\noindent\textbf{AB2 スタートアップ問題（最初の1ステップ目の特別処理）}

AB2 法は2時刻分の対流項履歴
$\bigl(\mathcal{C}^n,\,\mathcal{C}^{n-1}\bigr)$
を必要とするため，$t=0$ の最初のステップ（$n=0\to1$）では AB2 を適用できない．

\begin{center}
\renewcommand{\arraystretch}{1.2}
\begin{tabular}{lll}
\hline
ステップ & 対流項の扱い & 時間精度 \\
\hline
$n=0\to1$（初回）& 前進 Euler：$-\mathcal{C}(\bu^0,\bu^0)$ & $\Ord{\Delta t}$（スタートアップ） \\
$n=1\to2$ 以降 & AB2：$\tfrac{3}{2}\mathcal{C}^n - \tfrac{1}{2}\mathcal{C}^{n-1}$ & $\Ord{\Delta t^2}$ \\
\hline
\end{tabular}
\end{center}

初回ステップの前進 Euler による局所打ち切り誤差は $\Ord{\Delta t^2}$（1次精度スキームを1ステップのみ適用した標準値）であり，
以後の AB2 ステップへの伝播が\textbf{安定}であれば，全体積分精度 $\Ord{\Delta t^2}$ に影響しない．

\noindent\textbf{AB2 の零安定性による正当化：}
AB2 は線形多段スキームであり，特性多項式 $\rho(z) = z^2 - z$ の根は $z=0,\,z=1$ で，
いずれも $|z| \leq 1$ かつ $z=1$ は単根（Dahlquist の根条件を満足）．
これにより AB2 は\textbf{零安定（zero-stable）}であり，
$n=0$ での局所誤差 $e_1 = \Ord{\Delta t^2}$ は各後続ステップで高々定数倍（$\leq C$ で $C$ は問題に依存しない有界な定数）にしか増幅されない
（Dahlquist--Bjorck, \textit{Numerical Methods}, 1974 §5.3）．
結果として $N_t - 1$ ステップ後も $e_1$ の寄与は $\Ord{\Delta t^2}$ のまま保たれる．

\noindent\textbf{誤差積み上げの明示：}
$N_t = T/\Delta t$ ステップ積分で，
初回の局所誤差 $\Ord{\Delta t^2}$ の全体誤差への寄与は $\leq C\,\Ord{\Delta t^2}$（定数倍；$N_t$ 倍されない）；
残り $N_t - 1$ ステップの AB2 局所誤差 $\Ord{\Delta t^3}$ の合算は
$\Ord{(N_t-1)\Delta t^3} = \Ord{T\Delta t^2}$
であり，全体精度は $\Ord{\Delta t^2}$ に保たれる．

初回完了後，対流項 $\mathcal{C}^0 = \mathcal{C}(\bu^0,\bu^0)$ をバッファに保存し，
以降の AB2 ステップで利用する（追加メモリ：速度場1つ分，$2N_xN_y$ 実数値）．
\end{tcolorbox}

\subsubsection{Crank--Nicolson 法（粘性項）}

粘性項（2階空間微分を含む）は，純陽的処理では厳しい時間刻み制限（$\Delta t < O(h^2)$）を課す．
これを緩和するため，Crank--Nicolson（CN）法による\textbf{半陰的}処理を採用する：
\begin{equation}
  \frac{u^{n+1} - u^n}{\Delta t} = \frac{1}{2}\bigl[\mathcal{D}(u^n) + \mathcal{D}(u^{n+1})\bigr]
  \label{eq:crank_nicolson}
\end{equation}
ここで $\mathcal{D}$ は粘性拡散演算子．CN 法は時間方向に2次精度（$O(\Delta t^2)$）かつ
\textbf{無条件安定}（スペクトル半径の制約なし）であり，適度な $\Delta t$ での安定計算が可能となる．

\begin{tcolorbox}[warnbox, title={CN 粘性項のクロス微分：$u$ と $v$ の陰的結合と実装上の分離}]
\label{warn:cn_cross_derivative}

変粘性流体の粘性応力 $\nabla\cdot[\mu(\nabla\bu + \nabla^T\bu)]$ の $x$ 成分は：
\begin{align}
  \bigl[\nabla\cdot(\mu(\nabla\bu+\nabla^T\bu))\bigr]_x
  &= \underbrace{2\frac{\partial}{\partial x}\!\left(\mu\frac{\partial u}{\partial x}\right)
     + \frac{\partial}{\partial y}\!\left(\mu\frac{\partial u}{\partial y}\right)}_{\text{対角項 $\mathcal{D}_{\mathrm{d}}(u)$}}
   + \underbrace{\frac{\partial}{\partial y}\!\left(\mu\frac{\partial v}{\partial x}\right)}_{\text{クロス微分項 $\mathcal{D}_{\mathrm{c}}(v)$}}
  \label{eq:viscous_full_x}
\end{align}
クロス微分項 $\mathcal{D}_{\mathrm{c}}(v)$ を CN で $t^{n+1}$ に置くと $u^*$-$v^*$ の連立系（$2N^2$ 未知数）が生じ，ブロック Thomas 法の優位が失われる．\textbf{標準対処法：クロス項のみ時刻 $n$ で陽的評価し，対角項だけ CN で半陰的処理}する：
\begin{equation*}
  \frac{u^* - u^n}{\Delta t}
  = \frac{1}{2Re}\bigl[\mathcal{D}_{\mathrm{d}}(u^n) + \mathcal{D}_{\mathrm{d}}(u^*)\bigr]
  + \frac{1}{Re}\,\mathcal{D}_{\mathrm{c}}(v^n)
  + \bm{f}_x^{n+1}
\end{equation*}
これにより $u^*$, $v^*$ を独立にブロック Thomas 法で求解できる．

クロス微分項の1次陽的処理は $\Ord{\Delta t}$ 誤差を導入し，界面近傍（$|\bnabla\mu| \neq 0$）で全体精度を $\Ord{\Delta t}$ に劣化させる可能性がある．$\mu_l/\mu_g \gg 10$ の水/空気系では AB2 外挿 $[\frac{3}{2}\mathcal{D}_{\mathrm{c}}^n - \frac{1}{2}\mathcal{D}_{\mathrm{c}}^{n-1}]$ の適用を推奨する．
\end{tcolorbox}

\begin{tcolorbox}[warnbox, title={変粘性界面でのクロス微分陽的処理：粘性 CFL 条件の実質的復活}]
\label{warn:cross_cfl}

気液界面近傍では $\mu$ が数格子幅で $\mu_l/\mu_g \approx 10^2$ 変化する．クロス微分 $\mathcal{D}_{\mathrm{c}}(v)$ を陽的処理すると，界面上でのフォン・ノイマン安定性解析から局所 CFL 条件が復活する：
\begin{equation}
  \Delta t \le C_{\mathrm{cross}} \cdot \frac{h^2}{\Delta\mu / \rho}
  \label{eq:cross_cfl}
\end{equation}
ここで $C_{\mathrm{cross}} = O(1)$，$\Delta\mu = \mu_l - \mu_g$．$\mu_l/\mu_g = 100$ の系では等粘性系の陽的粘性 CFL より約 $1/100$ 倍厳しくなる可能性がある．

\begin{center}
\renewcommand{\arraystretch}{1.2}
\begin{tabular}{lll}
\hline
処理方法 & 安定条件 & 適用場面 \\
\hline
対角粘性項 $\mathcal{D}_{\mathrm{d}}$（CN 法） & 無条件安定 & 常に適用 \\
クロス微分項 $\mathcal{D}_{\mathrm{c}}$（陽的） & 式~\eqref{eq:cross_cfl} の制約あり & $\mu_l/\mu_g \lesssim 10$ では無害 \\
\hline
\end{tabular}
\end{center}

\textbf{シミュレーションが爆発した場合：}対流 CFL・毛管波制約（式~\eqref{eq:dt_adv},\eqref{eq:dt_sigma}）を確認後，なお爆発する場合はクロス微分項 $\mathcal{D}_{\mathrm{c}}$ が原因の可能性がある．$\Delta t$ を $1/10$ に減らして安定化するか確認する．
\end{tcolorbox}

% ─────────────────────────────────────────────────────────────────────────────
\subsection{CLS 再初期化：圧縮項の Dissipative CCD スキーム}
\label{sec:cls_compression}
% ─────────────────────────────────────────────────────────────────────────────

CLS 法の再初期化方程式（式 \eqref{eq:cls_reinit_iso}）
\begin{equation}
  \frac{\partial\psi}{\partial\tau}
  + \underbrace{\nabla\cdot\!\bigl(\psi(1-\psi)\hat{\bm{n}}\bigr)}_{\text{圧縮項（陽解法）}}
  = \underbrace{\nabla\cdot(\varepsilon\nabla\psi)}_{\text{拡散項（陰解法）}}
  \label{eq:cls_reinit_split}
\end{equation}
は\textbf{移流的圧縮項}（界面へ $\psi$ を引き戻す）と
\textbf{拡散項}（界面プロファイルを滑らかに保つ）の和として解釈できる．

\medskip
\noindent\textbf{陽解法・陰解法の使い分け：}
\begin{itemize}
  \item \textbf{圧縮項（陽解法）：}
        波速 $c(\psi) = \psi(1-\psi)|\hat{\bm{n}}|$ は $\psi = 0.5$ で最大値 $1/4$ をとり，
        バルク（$\psi \to 0$ または $1$）でゼロに退化する．
        §\ref{sec:advection_motivation} の Dissipative CCD 移流と同構成で陽的に処理する
        （CFL 制約：$\Delta\tau \leq 4h$ 程度）．
  \item \textbf{拡散項（陰解法）：}
        $\varepsilon\nabla^2\psi$ は $\Delta\tau < \varepsilon h^2$ という厳しい陽的 CFL 制約を持つ
        （典型値 $\varepsilon \approx h/2$ のとき $\Delta\tau < h^3/2$）．
        Crank--Nicolson（CN）による半陰的処理で無条件安定化し，
        圧縮項の CFL と同程度の大きな $\Delta\tau$ を許容する．
\end{itemize}

\medskip
\noindent\textbf{圧縮項の Dissipative CCD 離散化（陽解法）．}
$x$ 成分のフラックス場 $g_i = \psi_i(1-\psi_i)\hat{n}_{x,i}$ を定義し，
§\ref{sec:advection_motivation} と同一手順でフィルタ適用済みフラックス発散を計算する：
\begin{align}
  g'_i &= \left[\frac{\partial(\psi(1-\psi)\hat{n}_x)}{\partial x}\right]_i
  \quad\text{（CCD 求解）}
  \label{eq:comp_ccd} \\
  \tilde{G}_i &= g'_i + \varepsilon_d^{\mathrm{comp}}\!\left(g'_{i+1} - 2g'_i + g'_{i-1}\right)
  \quad\bigl(\varepsilon_d^{\mathrm{comp}} = \varepsilon_{d,\max} = 0.05\bigr)
  \label{eq:comp_filter}
\end{align}
$y$ 成分についても同様に処理し，合計フラックス発散
$\tilde{G} = \tilde{G}_x + \tilde{G}_y$ を得る．
法線ベクトル $\hat{\bm{n}} = \bnabla\psi / |\bnabla\psi|$ は
CCD + Dissipative フィルタ適用済みの 1 次微分場（§\ref{sec:dissipative_ccd}）
から評価する（ゼロ除算回避のため $|\bnabla\psi| + \delta$，$\delta = 10^{-8}$ を使用）．

\medskip
\noindent\textbf{拡散項の CCD-CN 陰解法（陰解法）．}
CCD Eq-II の2階微分行列（§\ref{sec:ccd_matrix}）を $M_2$（左辺係数，三重対角），
$B_2$（右辺係数，三重対角）と記す（$M_2\{\psi''_i\} = B_2\{h^{-2}\psi_i\}$）．
Crank--Nicolson の $\theta=1/2$ スキームにおける陰的拡散方程式は，
$M_2$ を両辺に乗じることで：
\begin{equation}
  \left(M_2 - \frac{\varepsilon\Delta\tau}{2h^2}\,B_2\right)\psi^{\tau+1}
  = \left(M_2 + \frac{\varepsilon\Delta\tau}{2h^2}\,B_2\right)\psi^{**}
  \label{eq:comp_cn_system}
\end{equation}
この線形系の係数行列 $M_2 - \frac{\varepsilon\Delta\tau}{2h^2}B_2$ は
\textbf{三重対角}（$M_2$，$B_2$ ともに三重対角のため）であり，
Thomas アルゴリズムで $O(N)$ に求解できる．ここで $\psi^{**}$ は
圧縮項の陽的処理後の中間値である．

\begin{tcolorbox}[algbox, title={CLS 再初期化の Dissipative CCD スキーム（1 仮想時間ステップ）}]
\label{alg:cls_reinit_dccd}

\textbf{入力：} $\{\psi_i^\tau\}$，$\varepsilon$（界面幅），$\varepsilon_d^{\mathrm{comp}}$，$\Delta\tau$

\begin{enumerate}
  \item \textbf{法線ベクトルの評価：}
        CCD + Dissipative フィルタ（§\ref{sec:dissipative_ccd}，$S(\psi)$ 制御）で
        $\{\bnabla\psi_i^\tau\}$ を計算；$\hat{n}_{x,i} = (\partial\psi/\partial x)_i / (|\bnabla\psi_i| + \delta)$．

  \item \textbf{圧縮フラックス（陽解法，Dissipative CCD）：}
        $g_i = \psi_i^\tau(1-\psi_i^\tau)\hat{n}_{x,i}$；
        CCD + 一様フィルタ（式~\eqref{eq:comp_ccd}--\eqref{eq:comp_filter}）で $\tilde{G}_{x,i}$ を計算；
        $y$ 方向も同様に $\tilde{G}_{y,i}$．

  \item \textbf{圧縮ステージ（陽的更新）：}
        \[
          \psi_i^{**} = \psi_i^\tau - \Delta\tau\!\left(\tilde{G}_{x,i} + \tilde{G}_{y,i}\right)
        \]
        クランプ：$\psi_i^{**} = \max(0,\min(1,\psi_i^{**}))$．

  \item \textbf{拡散ステージ（CN 陰解法，式~\eqref{eq:comp_cn_system}）：}
        右辺 $\bm{r} = (M_2 + \frac{\varepsilon\Delta\tau}{2h^2}B_2)\psi^{**}$ を計算；
        三重対角システム $(M_2 - \frac{\varepsilon\Delta\tau}{2h^2}B_2)\psi^{\tau+1} = \bm{r}$
        を Thomas アルゴリズムで求解（各方向独立に）．

  \item \textbf{クランプ：} $\psi_i^{\tau+1} = \max(0,\min(1,\psi_i^{\tau+1}))$．
\end{enumerate}

\textbf{仮想時間刻みの制約（圧縮項 CFL）：}
$\Delta\tau \leq C_\tau \cdot h$（$C_\tau = O(1)$ は波速の上界 $\max|\psi(1-\psi)|\hat{n}| \leq 1/4$ から）．
拡散項は CN 陰解法のため $\Delta\tau$ を制限しない（無条件安定）．
\end{tcolorbox}

% ─────────────────────────────────────────────────────────────────────────────
\subsection{安定性制約（時間刻みの決定）}
\label{sec:stability}
% ─────────────────────────────────────────────────────────────────────────────

時間刻み $\Delta t$ は，以下の3種類の制約のうち最も厳しい条件を選ぶ（安全率 $\text{CFL}_{\max} < 1$）：

\begin{tcolorbox}[mybox, title={時間刻みの制約（CN 法採用時：2制約）}]
\begin{align}
  &\text{\textbf{対流 CFL}：}\quad
  \Delta t_{\text{adv}} = \text{CFL}_{\max} \cdot \frac{1}{\max\!\left(\dfrac{|u|_{\max}}{\Delta x}+\dfrac{|v|_{\max}}{\Delta y}\right)}
  \label{eq:dt_adv} \\[6pt]
  &\text{\textbf{毛管波制約（2D）}：}\quad
  \Delta t_\sigma = \sqrt{\frac{(\rho_l + \rho_g)\min(\Delta x, \Delta y)^3}{2\pi\sigma}}
  \label{eq:dt_sigma}
\end{align}
\textbf{最終的な時間刻み：}$\Delta t = \min(\Delta t_{\text{adv}},\, \Delta t_\sigma)$

\medskip
\textbf{粘性 CFL（参考値）：}
\[
  \Delta t_{\text{visc}} = \frac{\min(\Delta x, \Delta y)^2}{4(\mu/\rho)_{\max}}
  \quad \text{（陽的粘性処理のみ適用；CN 法採用時は不要）}
\]
CN 法（無条件安定）を対角粘性項に採用しているため，等粘性流体では $\Delta t_{\text{visc}}$ は時間刻み決定に使用しない．
ただし，変粘性流体（気液二相流）ではクロス微分項 $\partial_y(\mu\partial_x v)$ が陽的処理されるため，
大密度比（$\mu_l/\mu_g \gg 1$）の場合には界面近傍で粘性類似の制約が実質的に復活しうる
（詳細は §\ref{warn:cn_cross_derivative} 末尾の安定性解析を参照）．
\end{tcolorbox}

式 \eqref{eq:dt_sigma} の毛管波制約は，界面上の微小振幅波の分散関係
$\omega^2 = \sigma k^3/(\rho_l+\rho_g)$ にナイキスト波数 $k_{\max}=\pi/\Delta x$
を代入し，群速度 CFL 条件 $c_g\,\Delta t \le \Delta x$ を適用して得られる
（導出の全手順と $\rho_l\gg\rho_g$ での物理的解釈は付録~\ref{app:capillary_cfl} 参照）．
$\Delta t_\sigma \propto \Delta x^{3/2}$ の依存性は分散関係 $\omega\propto k^{3/2}$ に由来し，
格子を半分にすると約 2.8 倍の制約強化をもたらす．

\noindent\textbf{注意：}
表面張力が支配的な問題（$We \ll 1$，小液滴，毛管現象）では，
毛管波（capillary wave）の位相速度が大きくなり CFL 条件が厳しくなる．
特に $\Delta t_\sigma \propto \Delta x^{3/2}$ の依存性（格子を半分にすると $\approx 2.8$ 倍制約が厳しくなる）
のため，高解像度計算では毛管波制約が支配的になることが多い．

\textbf{対策：}
\begin{itemize}
  \item 表面張力の半陰的処理（Implicit Surface Tension）による制約の緩和
  \item 問題によっては圧力 $p$ から $p + p_0$（静水圧 $p_0 = \rho g z$）を分離する
        動的圧力（kinematic pressure）分割で寄生流れを抑制しつつ $\Delta t$ を拡大
\end{itemize}
